\begin{landscape}
\begin{table}[p]
\caption{Primary and sensitivity estimates for wealth recovery and frailty progression.}\label{tab:s6_sensitivity}
\centering
\fontsize{9}{10.5}\selectfont
\setlength{\tabcolsep}{4pt}
\renewcommand{\arraystretch}{1.08}
\begin{tabularx}{\linewidth}{@{}Y L{3.8cm}rrrL{3.7cm}@{}}
\toprule
Analysis specification & Cohort & Participants, n & FI model observations, n & Adjusted difference in annual FI change, FI units/year & 95\% CI \\
\midrule
\multicolumn{6}{@{}l}{\textbf{Primary analysis}} \\
 & HRS & 5,114 & 23,411 & 0.0013 & 0.0006 to 0.0019 \\
 & ELSA & 699 & 2,584 & $-0.0014$ & $-0.0037$ to 0.0008 \\
 & KLoSA & 2,502 & 10,762 & 0.0006 & $-0.0006$ to 0.0017 \\
 & \textbf{Overall (random-effects)} & 3 cohorts & -- & 0.0005 & $-0.0026$ to 0.0035 \\
\addlinespace[2pt]
\multicolumn{6}{@{}l}{\textbf{Recovery threshold: 25\%}} \\
 & HRS & 5,114 & 23,411 & 0.0015 & 0.0009 to 0.0021 \\
 & ELSA & 699 & 2,584 & 0.0002 & $-0.0020$ to 0.0024 \\
 & KLoSA & 2,502 & 10,762 & 0.0009 & $-0.0002$ to 0.0020 \\
 & \textbf{Overall (random-effects)} & 3 cohorts & -- & 0.0012 & $-0.0001$ to 0.0026 \\
\addlinespace[2pt]
\multicolumn{6}{@{}l}{\textbf{Recovery threshold: 100\% (return to or above pre-shock wealth)}} \\
 & HRS & 5,114 & 23,411 & 0.0009 & 0.0002 to 0.0017 \\
 & ELSA & 699 & 2,584 & $-0.0028$ & $-0.0056$ to 0.0000 \\
 & KLoSA & 2,502 & 10,762 & 0.0018 & 0.0005 to 0.0030 \\
 & \textbf{Overall (random-effects)} & 3 cohorts & -- & 0.0002 & $-0.0054$ to 0.0059 \\
\addlinespace[2pt]
\multicolumn{6}{@{}l}{\textbf{Common-22 frailty index}} \\
 & HRS & 5,082 & 23,200 & 0.0008 & 0.0002 to 0.0014 \\
 & ELSA & 699 & 2,584 & $-0.0014$ & $-0.0033$ to 0.0005 \\
 & KLoSA & 2,502 & 10,762 & 0.0006 & $-0.0004$ to 0.0015 \\
 & \textbf{Overall (random-effects)} & 3 cohorts & -- & 0.0003 & $-0.0022$ to 0.0028 \\
\addlinespace[2pt]
\multicolumn{6}{@{}l}{\textbf{Landmark-weighted analysis}} \\
 & HRS & 5,097 & 23,361 & 0.0009 & $-0.0003$ to 0.0021 \\
 & ELSA & 676 & 2,508 & $-0.0014$ & $-0.0047$ to 0.0020 \\
 & KLoSA & 2,502 & 10,762 & 0.0007 & $-0.0008$ to 0.0021 \\
\addlinespace[2pt]
\multicolumn{6}{@{}l}{\textbf{Multiple-imputation analysis}} \\
 & HRS & 5,222 & 23,846 & 0.0012 & 0.0006 to 0.0019 \\
 & ELSA & 772 & 2,886 & $-0.0016$ & $-0.0037$ to 0.0006 \\
 & KLoSA & 2,507 & 10,781 & 0.0006 & $-0.0006$ to 0.0017 \\
 & \textbf{Overall (random-effects)} & 3 cohorts & -- & 0.0004 & $-0.0029$ to 0.0037 \\
\addlinespace[2pt]
\bottomrule
\end{tabularx}
\vspace{4pt}
\parbox{\linewidth}{\raggedright\fontsize{9}{11}\selectfont\textit{Notes:} Estimates are recovery-by-time interactions from cohort-specific longitudinal models and represent the adjusted difference in annual FI change (FI units/year) between substantial recovery and persistent/limited recovery. Negative estimates indicate slower annual FI increase in the substantial-recovery group. Primary models were unweighted. The landmark-weighted analysis used positive survey weights recorded at the landmark interview rescaled to a mean of one within cohort and participant-level CR0 sandwich standard errors. Cohort-specific weighted estimates were not pooled. Multiple-imputation estimates were combined within cohort using Rubin's rules. Participants and FI model observations may differ across specifications; FI model observations include landmark and eligible subsequent FI assessments. Recovery-threshold analyses reclassified substantial recovery at 25\% of shock-related wealth loss regained or at 100\% recovery to or above the pre-shock wealth level at the landmark. The Common-22 analysis used the 22 conceptually aligned frailty deficits across cohorts. Overall rows are secondary random-effects summaries of cohort-specific estimates and do not represent participant-level pooled models. The primary random-effects summary included three cohorts. Between-cohort heterogeneity was $\tau^2 = 9.40 \times 10^{-7}$ (FI units/year)$^2$ and $I^2 = 72.5\%$. FI, frailty index; MI, multiple imputation; HRS, Health and Retirement Study; ELSA, English Longitudinal Study of Ageing; KLoSA, Korean Longitudinal Study of Aging.}
\end{table}
\end{landscape}
